\documentclass[11pt]{article}                                    % Documento de artículo, tamaño de fuente 11pt
\usepackage[spanish]{babel}                                      % Idioma español
\usepackage[a4paper,margin=2.5cm]{geometry}                      % Márgenes
\usepackage{parskip,mlmodern,charter}                                     % Espaciado entre párrafos y fuente
\usepackage{mathtools,amssymb}                                   % Matemáticas avanzadas
\usepackage[output-decimal-marker={,}]{siunitx}                  % Unidades SI
\usepackage{graphicx,subcaption,booktabs,multirow}               % Imágenes y tablas
\usepackage{algorithmicx}                                        % Pseudocódigo
\usepackage{xcolor}                                              % Definición de colores
\usepackage{fancyhdr,lastpage}                                   % Encabezados y pies de página
\usepackage[numbers,square,sort]{natbib}                         % Bibliografía con estilo numérico
\usepackage[colorlinks=true, linkcolor=black, urlcolor=magenta, 
            citecolor=orange, anchorcolor=blue]{hyperref}        % Hipervínculos
\usepackage{microtype,cleveref}                                  % Mejoras tipográficas y referencias inteligentes

% cleverref: En vez de escribir la Figura \ref{fig:motor}, escribes la \cref{fig:motor}

\graphicspath{{./Imágenes/}} % Carpeta donde se encuentran las imágenes

\title{\textbf{Control de Accionamientos y Automatización de Grúa
Portacontenedores de Muelle tipo Pórtico} \\ \Large{Proyecto Global Integrador}}

\author{Aleo Rodrigo \\ Pucciarelli Nahuel \\ \\ AUTÓMATAS Y CONTROL DISCRETO 
\\ Universidad Nacional de Cuyo \\ Mendoza, Argentina}

\date{2026}

\begin{document}

 % Configuración de encabezados y pies de página
\pagestyle{fancy}
\fancyhf{}
\fancyfoot[R]{\thepage}
\renewcommand{\headrulewidth}{0pt} 
\urlstyle{same} 

\maketitle

\subsubsection*{Resumen}
Desarrollo e implementación de un Sistema Híbrido de Control y Protección con Operación Semi-Automática coordinada para una Grúa Portacontenedores Portuaria de Muelle tipo Pórtico, para carga y descarga de contenedores (containers) en grandes barcos de transporte marítimo.

El sistema híbrido de control y protección incluye 3 niveles jerárquicos:
\begin{itemize}
    \item Nivel 2 - Control Regulatorio: estados continuos, muestreados en tiempo discretizado (Sampled CTDS).
    \item Nivel 1 - Control Supervisor: autómata, estados discretos activados por eventos (DEDS).
    \item Nivel 0 - Protección: en espera; actúa en situación de emergencia, lleva el sistema a estado seguro.
\end{itemize}

\newpage
\tableofcontents
\newpage

\section{Desarrollo y modelado}
Este proyecto considera el control de accionamientos eléctricos y operación semi-automática coordinada de los dos movimientos principales continuos y simultáneos de la grúa en un plano vertical (2D): traslación horizontal del carro (Trolley) sobre las vigas principal y móvil (paralelo al eje $x$); e Izaje/descenso, o traslación vertical, de la carga (Hoist) bajo la posición del carro (paralelo al eje $y$).

Ambos movimientos son impulsados mediante cables de acero (wireropes) desde mecanismos rotativos con tambores, cajas de engranajes, frenos mecánicos y motores eléctricos, ubicados en una Casa de Máquinas (Machinery House) fija, con sus correspondientes accionamientos electrónicos de velocidad variable; dichos movimientos tienen \textit{restricciones de operación}: posiciones límite de recorrido, velocidades máximas y aceleraciones máximas. Están \textit{acoplados entre sí por la carga} (headblock/spreader con o sin container), que se balancea suspendida del carro que la traslada mediante los cables de izaje.

\subsubsection*{Hipótesis simplificativas}
Para el modelado dinámico de los movimientos de la grúa, se consideran las siguientes hipótesis simplificativas:
\begin{itemize}
    \item \textit{Estructura:} Se asume totalmente rígida; se desprecian las deformaciones y vibraciones de las vigas y el pórtico.
    \item \textit{Mecanismo de carro:} Se modela como una transmisión rígida accionada por un cable equivalente siempre en tensión, considerando su elasticidad longitudinal y amortiguamiento.
    \item \textit{Sistema de izaje:} Se simplifica a un sistema equivalente (un solo motor/tambor). El cable se considera de masa despreciable y comportamiento elástico solo a tracción (no soporta compresión).
    \item \textit{Carga:} Se modela como una masa puntual concentrada en la base del spreader, con 2 grados de libertad en el plano vertical (2D).
    \item \textit{Efectos externos:} Solo se consideran la gravedad y las fuerzas de contacto vertical (modelo elástico lineal). Se desprecia la resistencia aerodinámica frontal al viento de la carga y del carro en movimiento.
\end{itemize}

\subsection{Accionamiento y traslación del Carro}
\subsubsection{Dinámica del motor y freno (eje rápido)}
El motor genera torque, el freno se opone y la inercia resiste el cambio de velocidad. Esto se ve reflejado en la ecuación dinámica fundamental del eje del motor:
\begin{equation}
    J_{tm+tb} \, \frac{d\omega_{tm}}{dt} = T_{tm} + T_{tb} - b_{tm} \, \omega_{tm} - \frac{T_{td}}{i_t}
\end{equation}
donde $T_{tm}$ es el torque eléctico, $T_{tb}$ el del freno, y $\frac{T_{td}}{i_t}$ es el torque de carga reflejado desde el tambor. 
\subsubsection{Transmisión y tambor (eje lento)}
La caja reductora reduce la velocidad y aumenta el torque. El tambor convierte movimiento rotacional a lineal.
\begin{itemize}
    \item Relación de transmisión: $\omega_{td} = \frac{\omega_{tm}}{i_t}$
    \item Conversión a lineal (en el tambor): $v_{dt} = r_{td} \, \omega_{td}$. Nota: $v_dt$ es la \textit{velocidad lineal del cable en el tambor}, no la del carro.
\end{itemize}
\subsubsection{Elasticidad del cable (acoplamiento)}
El cable no es rígido, actúa como resorte ($K_{tw}$) y amortiguador ($b_{tw}$) entre la posicion demandada por el tambor ($x_{td}$) y la posición real del carro ($x_t$). La fuerza de tracción del cable es:
\begin{equation}
    F_{tw} = K_{tw} \, (x_{td} - x_t) + b_{tw} \, (v_{td} - v_t)
\end{equation}
\subsubsection{Dinámica del carro (masa traslacional)}
Finalmente, se aplica la segunda ley de Newton al carro:
\begin{equation}
    M_t \, \frac{dv_t}{dt} = F_{tw} - b_t \, v_t + F_{coupling}
\end{equation}
donde $F_{coupling} = 2\,F_{hw} \, \sin\theta_l$.
\subsection{Accionamiento y traslación del Izaje}

\end{document}